\documentclass[11pt]{article}
\usepackage[margin=0.85in]{geometry}
\usepackage[T1]{fontenc}
\usepackage[utf8]{inputenc}
\usepackage{amsmath,amssymb}
\usepackage{array}
\usepackage{longtable}
\usepackage{booktabs}
\usepackage{enumitem}
\setlist[itemize]{topsep=3pt,itemsep=2pt,parsep=0pt}

\title{Stage 03 Full Theory-to-Code Map v1}
\author{Reaction-boundary volatility surface project}
\date{5 July 2026}

\begin{document}
\maketitle

\section*{Artefact Status}
Full Stage 3 theory-to-code map. This is still a planning artefact. No
code, figures, tables, parameter defaults, numerical diagnostics, or
mathematical verification have been produced.

\section*{Source Numbering Note}
The user refers to the main surface-generating result as Eq.~18 and the
response-function asymptotic cases as Eq.~22. In the frozen source, the
stable labels are:
\begin{itemize}
  \item Main-result closure: \texttt{eq:main-operational-result} in the
  main results section, repeated as \texttt{eq:operational-result} in the
  derivation.
  \item Response-function integral: \texttt{eq:F-definition}.
  \item Two response-function regimes: \texttt{eq:F-asymp}.
  \item Deterministic-clock projection: \texttt{eq:deterministic-localvol}.
  \item Operational kernel notation: \texttt{eq:Xi}.
\end{itemize}
Later code and run notes should cite the source labels as well as any
paper equation numbers.

\section*{Core Computational Route}
The first implementation should not simulate the signed-order-flow path or
the Abel response path directly. It should implement the reduced-form
surface engine:
\[
 a_u^{(\Delta)}(S,u)
 \simeq
 \left[\frac{\mathcal A_{\rm eff}(S,u)}{\mathcal L_u(S,u)^2}\right]
 \Delta^{-\gamma(S,u)}
 \mathcal F_{\gamma(S,u)}(\nu_u(S,u)\Delta),
\]
then apply the deterministic-clock projection
\[
 \sigma_{\rm loc}^2(S,t)=\alpha(t)\Xi(S,U(t);\Delta).
\]

\begin{longtable}{p{0.15\linewidth}p{0.21\linewidth}p{0.25\linewidth}p{0.27\linewidth}}
\toprule
Source label & Theory object & Computational object & Required gate before code \\
\midrule
\texttt{eq:main-operational-result}, \texttt{eq:operational-result} &
Operational variance closure \(a_u^{(\Delta)}(S,u)\) &
Function target: evaluate operational variance over a state/maturity or state/time grid &
Stage 5 must choose state axis, parameter profiles, and response-regime branch. \\
\midrule
\texttt{eq:Xi} &
Operational kernel \(\Xi(S,u;\Delta)\) &
Alias or wrapper for the operational closure used in clock projection &
Must remain an order-book object until the clock multiplier is applied. \\
\midrule
\texttt{eq:F-asymp} &
Low-\(z\) response case \( \mathcal F_\gamma(z)\sim \tan(\pi\gamma/2)/[2(1-\gamma)]\) &
Response regime A: small \(z=\nu\Delta\), constant in \(z\) for fixed \(\gamma\) &
Need admissible \(\gamma\) range away from \(1\); record low-\(z\) condition. \\
\midrule
\texttt{eq:F-asymp} &
High-\(z\) response case \( \mathcal F_\gamma(z)\sim K_\gamma^{(0)}z^{\gamma-1}\) &
Response regime B: large \(z=\nu\Delta\), resilience-scale decay for \(0<\gamma<1\) &
Need implementation of \(K_\gamma^{(0)}\) and admissible \(\gamma\) range. \\
\midrule
\texttt{eq:Kgamma} &
High-\(z\) prefactor \(K_\gamma^{(0)}\) &
Helper function target if high-\(z\) branch is implemented &
Requires gamma-function library choice and endpoint safeguards. \\
\midrule
\texttt{eq:deterministic-localvol} &
Calendar-time local variance \(\sigma_{\rm loc}^2(S,t)\) &
Deterministic clock projection function &
Need \(\alpha(t)\) and \(U(t)\) convention. Use deterministic clock only. \\
\midrule
\texttt{eq:homogeneous-result} &
Homogeneous benchmark \(\sigma_{\rm eff}^2\) &
Scalar diagnostic and baseline sanity check &
Use for baseline comparison, not as the full surface if state dependence is needed. \\
\midrule
\texttt{eq:linear-book} and explanatory text after \texttt{eq:operational-result} &
Liquidity slope suppresses motion through \(1/\mathcal L(S)^2\) &
Primary candidate for stock-price dependence: declare \(\mathcal L(S)\) profile &
Stage 5 must specify the profile and label it as a modelling choice. \\
\midrule
\texttt{eq:Aeff} &
\(\mathcal A_{\rm eff}=A_m/D_u\) &
Effective signed-forcing intensity parameter or profile &
Decide constant versus state/time dependence. \\
\midrule
\texttt{eq:coarse-rate} &
Finite operational scale \(\Delta\) &
Configuration parameter; not calendar maturity &
Record units and grid; do not confuse with \(T\). \\
\bottomrule
\end{longtable}

\section*{First Surface Convention To Decide Later}
The first surface should be generated over one of the following:
\begin{itemize}
  \item \(S=\log x\) and maturity \(T\);
  \item stock price \(x\) and maturity \(T\), with \(S=\log x\) internally;
  \item log-moneyness and maturity \(T\).
\end{itemize}
The map recommends stock price or log-moneyness for reader-facing figures,
while keeping \(S=\log x\) internally because the source theory uses \(S\)
as the boundary/log-price state.

\section*{Refactored Stock-Dependence Strategy}
Stock dependence should be generated by explicit structural profiles, not
by hiding all shape in a single surface. The first numerical design should
therefore prepare three simulation profiles:
\begin{enumerate}
  \item Liquidity-only: vary \(\mathcal L(S)\), hold
  \(\mathcal A_{\rm eff}\), \(\nu\), \(\gamma\), and \(\alpha\) fixed.
  This isolates the inverse-square liquidity channel.
  \item Resilience-response-only: vary \(\nu(S)\), hold
  \(\mathcal A_{\rm eff}\), \(\mathcal L\), \(\gamma\), and \(\alpha\)
  fixed. This isolates the \(z(S)=\nu(S)\Delta\) channel in
  \(\mathcal F_\gamma(z)\), especially in the high-\(z\) regime.
  \item Combined structural: vary \(\mathcal L(S)\), \(\nu(S)\), and
  optionally \(\mathcal A_{\rm eff}(S)\) together to produce a stylised
  stress or skew surface. This profile must be labelled as simulation
  phenomenology rather than calibration.
\end{enumerate}

\section*{Mandatory Later Outputs}
\begin{itemize}
  \item Parameter-range table, with raw CSV and standalone LaTeX table.
  \item Response-regime table for the two \(\mathcal F_\gamma\) cases.
  \item Matched surface-and-contour output for the low-\(z\) response case.
  \item Matched surface-and-contour output for the high-\(z\) response case.
\end{itemize}

\section*{Deferred Objects}
The spectral cumulant, finite-cutoff response, pricing PDE, stochastic or
non-unique clocks, empirical calibration, and direct signed-order-flow path
simulation are deferred. They should not enter the first implementation
unless a later stage explicitly revises scope.

\end{document}
